\documentclass[11pt]{article}

% Packages
\usepackage[utf8]{inputenc}
\usepackage[T1]{fontenc}
\usepackage{amsmath,amssymb,amsfonts}
\usepackage{graphicx}
\usepackage{booktabs}
\usepackage{hyperref}
\usepackage[margin=1in]{geometry}
\usepackage{authblk}
\usepackage{caption}
\usepackage{subcaption}
\usepackage{algorithm}
\usepackage{algpseudocode}
\usepackage{xcolor}
\usepackage{natbib}

% Title
\title{MPKNet: A Lateral Geniculate Nucleus-Inspired Architecture for Parameter-Efficient Visual Processing}

\author[1]{D.J. Lougen}
\affil[1]{University of Toronto \\ \texttt{d.lougen@mail.utoronto.ca}}

\date{}

\begin{document}

\maketitle

\begin{abstract}
We present MPKNet, a biologically-inspired convolutional neural network that models the parallel processing streams of the Lateral Geniculate Nucleus (LGN). Unlike conventional ``bio-inspired'' approaches that borrow surface-level features such as Gabor filters, MPKNet directly implements the laminar organization observed across mammalian visual systems---Magnocellular (M), Parvocellular (P), and Koniocellular (K) pathways---as separate processing streams with distinct spatial sampling densities. The architecture achieves 40.6\% accuracy on TinyImageNet-200 with only 0.21M parameters, matching ResNet18 (41.5\%) with 52$\times$ fewer parameters. On medical imaging (Kvasir-v2), it achieves 89.2\% with no pretraining or augmentation. Notably, data augmentation \textit{degrades} performance (40.6\% $\rightarrow$ 24.1\% on TinyImageNet), suggesting the architecture provides intrinsic invariances that external augmentation disrupts. This work was conducted independently during summer 2025 using personal hardware (DGX Spark, MacBook M3 Max), demonstrating that meaningful architectural research is possible without large compute clusters.
\end{abstract}

\section{Introduction}

The mammalian visual system processes information through parallel pathways before integration in primary visual cortex (V1). The Lateral Geniculate Nucleus (LGN), positioned between the retina and V1, contains three distinct cell populations that have been conserved across 200 million years of mammalian evolution \citep{solomon2021parallel}:

\begin{itemize}
    \item \textbf{Magnocellular (M) cells}: Large receptive fields, fast temporal responses, sensitive to luminance contrast and motion
    \item \textbf{Parvocellular (P) cells}: Small receptive fields, sustained responses, high spatial acuity and color sensitivity
    \item \textbf{Koniocellular (K) cells}: Sparse, heterogeneous population involved in cross-stream modulation
\end{itemize}

This parallel organization is remarkably consistent from tree shrews to humans \citep{conley1984tree}, suggesting evolutionary optimization for efficient visual processing. Yet most ``bio-inspired'' neural networks borrow only surface-level features (Gabor-like filters, center-surround operations) without implementing the fundamental architectural principle: \textit{parallel specialist streams with late fusion}.

We hypothesize that the \textit{pattern} of information routing---which neurons multiply with which---encodes useful computational priors. Standard CNNs allow every neuron to interact with every other in subsequent layers. MPKNet restricts this: M talks to M, P talks to P, K modulates but does not mix features. The math is identical; the wiring diagram differs.

This work was inspired by Yamins et al. \citep{yamins2014performance}, who demonstrated that performance-optimized hierarchical models naturally develop representations similar to primate visual cortex. We extend this by asking: if we \textit{start} with biologically-accurate connectivity constraints, what computational properties emerge?

\section{Related Work}

\paragraph{Bio-Inspired Vision Models.} Classical bio-inspired approaches include HMAX \citep{riesenhuber1999hierarchical}, which models simple and complex cells, and Gabor filter banks inspired by V1 receptive fields. More recent work explores spiking neural networks \citep{tavanaei2019deep} and predictive coding frameworks \citep{rao1999predictive}. However, these typically focus on individual mechanisms rather than the global organization of parallel visual streams.

\paragraph{Efficient Architectures.} MobileNet \citep{howard2017mobilenets}, EfficientNet \citep{tan2019efficientnet}, and SqueezeNet \citep{iandola2016squeezenet} achieve efficiency through depthwise separable convolutions, neural architecture search, and compression. These are engineering solutions without biological motivation. We show that biological constraints can achieve comparable efficiency while exhibiting qualitatively different training dynamics.

\paragraph{LGN Modeling.} Prior work on LGN modeling has focused on single-cell responses \citep{carandini2005we} or M/P pathway differences in temporal processing \citep{derrington1984spatial}. To our knowledge, MPKNet is the first to implement the full parallel-stream organization as a deep learning architecture.

\section{Architecture}

\subsection{Core Principle: Stride-Based Pathway Differentiation}

The key insight of MPKNet V4 (``MPKx'') is that M, P, and K pathways differ primarily in their \textit{spatial sampling density}, not kernel shape. Biologically, M cells tile the retina more sparsely than P cells. We implement this through stride:

\begin{itemize}
    \item \textbf{P pathway}: stride=1 (dense sampling, fine detail)
    \item \textbf{K pathway}: stride=2 (intermediate density)
    \item \textbf{M pathway}: stride=3 (sparse sampling, global gist)
\end{itemize}

All pathways use identical 3$\times$3 kernels. This eliminates mid-network pooling---pathways maintain their natural resolutions until V1 fusion, with only Global Average Pooling (GAP) at the classification head.

\subsection{Binocular Processing}

The LGN exhibits eye-specific organization: layers 1, 4, 6 receive contralateral input while layers 2, 3, 5 receive ipsilateral input. MPKNet implements this through:

\begin{enumerate}
    \item \textbf{Stereo disparity simulation}: Single images are split into left/right views with horizontal shifts (range $\pm$2 pixels), providing depth cues during training
    \item \textbf{Eye segregation}: Left and right pathways remain separate through LGN processing
    \item \textbf{V1 fusion}: Eye-specific streams are combined only at the final stage, mimicking binocular integration in V1
\end{enumerate}

\subsection{Retinal Preprocessing (BinocularPreMPK)}

Before pathway processing, input images undergo center-surround filtering to simulate retinal ganglion cell responses:

\begin{align}
    P &= I - G_\sigma * I \quad \text{(high-pass, detail)} \\
    M &= G_\sigma * \bar{I} \quad \text{(low-pass luminance, gist)}
\end{align}

where $G_\sigma$ is a Gaussian blur kernel and $\bar{I}$ is the luminance (mean across color channels). This separates the input into P-like (edges, detail) and M-like (global luminance, motion-relevant) representations.

\subsection{K-Gating Mechanism}

The Koniocellular pathway generates channel attention gates that modulate both P and M streams:

\begin{align}
    g_M &= \sigma(W_M \cdot \text{GAP}(K)) \\
    g_P &= \sigma(W_P \cdot \text{GAP}(K)) \\
    M' &= M \odot g_M, \quad P' = P \odot g_P
\end{align}

This implements context-aware gain control, where K acts as a relay that adjusts the relative importance of detail vs. global information based on scene content.

\subsection{V1 Fusion}

At the output, pathways are pooled to matching spatial dimensions and concatenated:

\begin{equation}
    z = \text{Conv}_{1\times1}([M_L \| M_R \| P_L \| P_R])
\end{equation}

followed by GAP and a linear classifier. This late fusion preserves pathway-specific representations until the final integration stage.

\subsection{Model Summary}

\begin{table}[h]
\centering
\begin{tabular}{ll}
\toprule
\textbf{Metric} & \textbf{Value} \\
\midrule
Parameters & 0.21--0.22M \\
Model size & 0.89 MB \\
Embedding dimension & 96 floats (384 bytes) \\
FLOPs & 142--280M (task dependent) \\
\bottomrule
\end{tabular}
\caption{MPKx architecture summary}
\end{table}

\section{Experiments}

All experiments were conducted on a DGX Spark or MacBook M3 Max, using personal hardware without university compute resources. Models were trained from scratch with no pretraining.

\subsection{TinyImageNet-200}

TinyImageNet contains 200 classes with 64$\times$64 images. Results are shown in Table~\ref{tab:tinyimagenet}.

\begin{table}[h]
\centering
\begin{tabular}{lcccc}
\toprule
\textbf{Model} & \textbf{Params} & \textbf{Test Acc} & \textbf{Acc/Param} & \textbf{Notes} \\
\midrule
\textbf{MPKx} & \textbf{0.21M} & \textbf{40.6\%} & \textbf{193} & No augmentation \\
MPKx (with aug) & 0.21M & 24.1\% & 115 & Augmentation hurts \\
ResNet18 & 11M & 41.5\% & 3.8 & 52$\times$ more params \\
MobileNetV2 & 3.4M & 33.1\% & 9.7 & 16$\times$ more params \\
EfficientNet-B0 & -- & 36.9\% & -- & \\
\bottomrule
\end{tabular}
\caption{TinyImageNet-200 results. MPKx matches ResNet18 with 52$\times$ fewer parameters. Notably, augmentation \textit{degrades} performance by 16.5 points.}
\label{tab:tinyimagenet}
\end{table}

\paragraph{Augmentation Sensitivity.} Standard data augmentation (random crops, flips, color jitter) dramatically hurts MPKx performance: 40.6\% $\rightarrow$ 24.1\%. This is consistent across datasets (CIFAR-100: 52.8\% without augmentation vs. 46.0\% with). We hypothesize that the parallel M/P/K architecture provides intrinsic invariances that conflict with augmentation-learned invariances.

\paragraph{Training Dynamics.} MPKx exhibits an unusually tight train/test gap ($\sim$3.4\%) throughout training, despite no augmentation. The model does not memorize the training set in the way typical CNNs do.

\subsection{Medical Imaging: Kvasir-v2}

Kvasir-v2 contains 8 classes of gastrointestinal endoscopy images. Results are shown in Table~\ref{tab:kvasir}.

\begin{table}[h]
\centering
\begin{tabular}{lccc}
\toprule
\textbf{Model} & \textbf{Params} & \textbf{Accuracy} & \textbf{Acc/Param} \\
\midrule
\textbf{MPKx} & \textbf{0.21M} & \textbf{89.2\%} & \textbf{425} \\
ConvMixer & 0.59M & 92.5\% & 157 \\
SqueezeNet & 1.25M & 85.6\% & 68 \\
MobileNetV3-Small & 2.5M & 92.5\% & 37 \\
DenseNet201 & 19.2M & 94.5\% & 5 \\
\bottomrule
\end{tabular}
\caption{Kvasir-v2 results. MPKx achieves 89.2\% with no pretraining or augmentation, within 3.3 points of ConvMixer at 3$\times$ fewer parameters.}
\label{tab:kvasir}
\end{table}

\subsection{Cross-Dataset Summary}

\begin{table}[h]
\centering
\begin{tabular}{lccccc}
\toprule
\textbf{Dataset} & \textbf{Params} & \textbf{FLOPs} & \textbf{Test Acc} & \textbf{Acc/Param} & \textbf{Train/Test Gap} \\
\midrule
TinyImageNet & 0.21M & 142M & 40.6\% & 193 & 3.5\% \\
CIFAR-100 & 0.22M & 281M & 58.8\% & 267 & 11\% \\
Kvasir-v2 & 0.21M & 280M & 89.2\% & 425 & 8\% \\
STL-10 & 0.21M & 280M & 71.7\% & 341 & 12\% \\
\bottomrule
\end{tabular}
\caption{MPKx cross-dataset performance. All results with no pretraining and no augmentation.}
\label{tab:summary}
\end{table}

\subsection{Ablation Study (STL-10)}

To validate that each pathway contributes meaningfully, we conducted ablations on STL-10 (Table~\ref{tab:ablation}).

\begin{table}[h]
\centering
\begin{tabular}{lcccc}
\toprule
\textbf{Ablation} & \textbf{Test Acc} & \textbf{Train Acc} & \textbf{Gap} & \textbf{$\Delta$ from Full} \\
\midrule
Full model & 71.0\% & 85\% & 14pt & -- \\
No K-gating & 71.1\% & 80\% & 9pt & +0.1pt \\
No M pathway & 70.0\% & 81\% & 11pt & $-$1.0pt \\
No P pathway & 62.8\% & 71\% & 8pt & $-$8.2pt \\
\bottomrule
\end{tabular}
\caption{Pathway ablation study on STL-10.}
\label{tab:ablation}
\end{table}

\paragraph{Interpretation.}
\begin{itemize}
    \item \textbf{P is load-bearing} ($-$8.2pt): Removing the Parvocellular pathway severely impairs classification, consistent with P-cells' role in fine spatial discrimination.
    \item \textbf{M carries generalizable information} ($-$1.0pt): The Magnocellular pathway contributes real signal that transfers to test data, suggesting useful global context.
    \item \textbf{K-gating adds capacity that doesn't generalize on static images} (+0.1pt test, but +5pt train): K-gating increases training accuracy without improving test performance. This suggests K's modulatory role is optimized for dynamic/temporal contexts---static image benchmarks may underestimate K's contribution to real visual processing.
\end{itemize}

\subsection{Prototype-Based Retrieval}

MPKx embeddings support nearest-prototype classification without retraining:

\begin{table}[h]
\centering
\begin{tabular}{lcc}
\toprule
\textbf{Evaluation} & \textbf{Prototype Acc} & \textbf{Linear Acc} \\
\midrule
Train (held-out 20\%) & 71.2\% & 88.7\% \\
Test set & 49.8\% & 55.8\% \\
\bottomrule
\end{tabular}
\caption{CIFAR-100 prototype-based retrieval. The 71\% prototype accuracy shows MPKx learns retrieval-ready embeddings.}
\end{table}

The 96-float embeddings are 5--20$\times$ more compact than CLIP/ResNet (512--2048 floats), making them suitable for edge deployment and image retrieval applications.

\section{Discussion}

\subsection{Why Does Augmentation Hurt?}

The consistent degradation with augmentation (TinyImageNet: $-$16.5pt; CIFAR-100: $-$6.8pt) is striking. We speculate that:

\begin{enumerate}
    \item The M/P/K architecture provides intrinsic invariances through its parallel processing structure
    \item External augmentation creates conflicting training signals that disrupt these natural invariances
    \item The biological visual system similarly relies on architectural invariances rather than learned ones
\end{enumerate}

This requires further investigation but suggests a fundamentally different approach to model regularization.

\subsection{Biological Implications}

The K-gating ablation result is consistent with known K-cell function: K-cells are involved in slow temporal tasks (circadian rhythm, ambient light level) rather than rapid spatial discrimination. Static image benchmarks naturally underestimate their contribution.

\subsection{Compute Democratization}

All experiments were run on personal hardware (DGX Spark, MacBook M3 Max) during summer 2025. The architecture was designed specifically to enable meaningful research without compute cluster access. With 0.21M parameters and $\sim$280M FLOPs, training completes in reasonable time on consumer hardware.

\section{Limitations and Future Work}

\paragraph{What This Work Ignores.} Several biological features are deliberately omitted:
\begin{itemize}
    \item \textbf{Temporal dynamics}: M cells respond transiently, P cells sustained---not modeled
    \item \textbf{Recurrence}: Massive feedback from V1 to LGN is not implemented
    \item \textbf{Foveation}: Varying resolution across visual field not modeled
    \item \textbf{Color opponency}: True red-green opponent channels not implemented
\end{itemize}

\paragraph{Future Directions.}
\begin{itemize}
    \item \textbf{V1 modeling}: Adding orientation columns and horizontal connections
    \item \textbf{Retinotectal pathway}: Superior colliculus-like saccade generation
    \item \textbf{Thalamo-cortical loops}: Testing whether attention emerges from architecture
    \item \textbf{Parallel execution}: M/P/K pathways are independent until fusion and could theoretically execute on separate threads
\end{itemize}

\section{Conclusion}

MPKNet demonstrates that biological organizational principles---parallel specialist streams with late fusion---translate to concrete computational advantages. The architecture matches ResNet18 on TinyImageNet with 52$\times$ fewer parameters, and exhibits qualitatively different training dynamics including augmentation insensitivity and tight train/test gaps.

More broadly, this work suggests that \textit{where} computation happens (the wiring diagram) may be as important as \textit{what} computation happens (the math). Evolution optimized visual processing for 200 million years; perhaps there is more to learn from its solutions than Gabor filters.

\section*{Acknowledgements}

Thanks to Paul Dassonville (University of Oregon) for introducing me to LGN cell types, and to Jay Pratt (University of Toronto) for ongoing collaboration on koniocellular function.

\bibliographystyle{plainnat}
\begin{thebibliography}{20}

\bibitem[Carandini \& Heeger(2005)]{carandini2005we}
Carandini, M., \& Heeger, D.~J. (2005).
\newblock Normalization as a canonical neural computation.
\newblock \textit{Nature Reviews Neuroscience}, 13(1), 51--62.

\bibitem[Conley et al.(1984)]{conley1984tree}
Conley, M., Fitzpatrick, D., \& Diamond, I.~T. (1984).
\newblock The laminar organization of the lateral geniculate body and the striate cortex in the tree shrew.
\newblock \textit{Journal of Neuroscience}, 4(1), 171--197.

\bibitem[Derrington \& Lennie(1984)]{derrington1984spatial}
Derrington, A.~M., \& Lennie, P. (1984).
\newblock Spatial and temporal contrast sensitivities of neurones in lateral geniculate nucleus of macaque.
\newblock \textit{The Journal of Physiology}, 357(1), 219--240.

\bibitem[Howard et al.(2017)]{howard2017mobilenets}
Howard, A.~G., Zhu, M., Chen, B., et al. (2017).
\newblock MobileNets: Efficient convolutional neural networks for mobile vision applications.
\newblock \textit{arXiv preprint arXiv:1704.04861}.

\bibitem[Iandola et al.(2016)]{iandola2016squeezenet}
Iandola, F.~N., Han, S., Moskewicz, M.~W., et al. (2016).
\newblock SqueezeNet: AlexNet-level accuracy with 50x fewer parameters and $<$0.5MB model size.
\newblock \textit{arXiv preprint arXiv:1602.07360}.

\bibitem[Rao \& Ballard(1999)]{rao1999predictive}
Rao, R.~P., \& Ballard, D.~H. (1999).
\newblock Predictive coding in the visual cortex: a functional interpretation of some extra-classical receptive-field effects.
\newblock \textit{Nature Neuroscience}, 2(1), 79--87.

\bibitem[Riesenhuber \& Poggio(1999)]{riesenhuber1999hierarchical}
Riesenhuber, M., \& Poggio, T. (1999).
\newblock Hierarchical models of object recognition in cortex.
\newblock \textit{Nature Neuroscience}, 2(11), 1019--1025.

\bibitem[Solomon(2021)]{solomon2021parallel}
Solomon, S.~G. (2021).
\newblock Parallel processing in the visual system.
\newblock \textit{Current Biology}, 31(11), R640--R647.

\bibitem[Tan \& Le(2019)]{tan2019efficientnet}
Tan, M., \& Le, Q. (2019).
\newblock EfficientNet: Rethinking model scaling for convolutional neural networks.
\newblock \textit{ICML}.

\bibitem[Tavanaei et al.(2019)]{tavanaei2019deep}
Tavanaei, A., Ghodrati, M., Kheradpisheh, S.~R., et al. (2019).
\newblock Deep learning in spiking neural networks.
\newblock \textit{Neural Networks}, 111, 47--63.

\bibitem[Yamins et al.(2014)]{yamins2014performance}
Yamins, D.~L., Hong, H., Cadieu, C.~F., et al. (2014).
\newblock Performance-optimized hierarchical models predict neural responses in higher visual cortex.
\newblock \textit{Proceedings of the National Academy of Sciences}, 111(23), 8619--8624.

\end{thebibliography}

\end{document}
